\newpage
\subsection{Interesting Problems}

\begin{exm}\label{exm:RetainedCustomers}
Retained customers

In the following, when we do the \code{LAG}, then SQL knows
how to behave, i.e. \code{prev\_active\_dt} is the previous day that a user was active, not just simply 1-day-before-current-day:
if user 1 is active on Jan 2, 2026 but not on Jan 1, 2026, in the new column \code{prev\_active\_date} the row that contains date = Jan 2, 2026, the \code{prev\_active\_date} = NULL.
Said differently, the \code{LAG()} function is not just subtracting current date by 1.

\begin{lstlisting}[language=SQL, linewidth=1.5\textwidth]
-- 1-day retention. 
WITH t AS (SELECT date_, ID, LAG(date_) OVER (PARTITION BY ID ORDER BY date_) AS prev_active_dt
           FROM activity)
SELECT date_, COUNT(DISTINCT ID) AS active_users,
 COUNT(DISTINCT CASE WHEN prev_active_dt=DATE(date_,'-1 day') THEN ID END) AS retained_u,
 COUNT(DISTINCT CASE WHEN prev_active_dt=DATE(date_,'-1 day') THEN ID END)*1.0/COUNT(DISTINCT ID) AS retention
FROM t
GROUP BY date_
ORDER BY date_;
\end{lstlisting}

\begin{lstlisting}[language=SQL, linewidth=1.5\textwidth]
-- Cohort retention: D7 retention : returned within 7 days
-- We can just change 7 to 1 to get above.
WITH t AS (SELECT date_, ID, LAG(date_) OVER (PARTITION BY ID ORDER BY date_) AS prev_dt
           FROM activity)
SELECT date_, COUNT(DISTINCT ID) AS active_users,
       COUNT(DISTINCT CASE WHEN prev_dt IS NOT NULL AND julianday(date_)-julianday(prev_dt) <= 7
               THEN ID END) AS retained_7d,
       COUNT(DISTINCT CASE WHEN prev_dt IS NOT NULL AND julianday(date_)-julianday(prev_dt) <= 7
               THEN ID END) * 1.0 / COUNT(DISTINCT ID) AS retention
FROM t
GROUP BY date_
ORDER BY date_;

-- in MySQL DATEDIFF(date_, prev_dt) <= 7
\end{lstlisting}
\end{exm}

\begin{exm}\label{exm:2059}
({\bf  \color{red} ID 2059: \href{https://platform.stratascratch.com/coding/2059-player-with-longest-streak?code_type=3}{Player with Longest Streak}})\sidenote{This was difficult. Needs creativity \idea}
\label{ex:playersLongestStreak}
\\
You are given a table of tennis players and their matches that they could either win (W) or lose (L). 
Find the longest streak of wins. A streak is a set of consecutive won matches of one player. 
The streak ends once a player loses their next match.

For this question, disregard edge cases such as: players who never lose, streaks that start before the first loss, and streaks that continue after the final match.

First, let us see how the table looks like below. Ignore the column called Loss Group; it is not
in the provided table. That column is added in the first step of our code. 
Each time a player loses, we increment a counter.
That naturally separates win streaks into groups.
Then, we count consecutive W's within each \code{loss\_group}.
Then, take the maximum.
\begin{table}[h]
\centering
\rowcolors{1}{aliceblue}{white}
\begin{tabular}{c|c|c|c}
\rowcolor{shadecolor} Player ID & Match Date & Result & Loss Group \\
\hline
401 & 2021-05-04 & W & 0 \\
401 & 2021-05-09 & L & 1 \\
401 & 2021-05-16 & L & 2 \\
401 & 2021-05-18 & W & 2 \\
401 & 2021-05-22 & L & 3 \\
401 & 2021-06-15 & L & 4 \\
401 & 2021-06-16 & W & 4 \\
401 & 2021-06-18 & W & 4 \\
401 & 2021-07-06 & L & 5 \\
401 & 2021-07-13 & L & 6 \\
\hline
\end{tabular}
\caption{Match Results for Player 401}
\end{table}

\vspace{.2in}
The second table, \code{win\_streaks}, looks like
\begin{table}[h]
\centering
\rowcolors{1}{aliceblue}{white}
\begin{tabular}{c|c|c}
\rowcolor{shadecolor} Player ID & loss group & streak length \\
\hline
401 & 0 & 1\\
401 & 2 & 1\\
401 & 4 & 2\\
\hline
\end{tabular}
\caption{Win streaks}
\end{table}

\noindent {\bf Reminder}
The following line acts like a running window, but, since no window
is specified it, by default, does running total:
\begin{lstlisting}[language=SQL, linewidth=1.4\textwidth]
SUM(CASE WHEN match_result = 'L' THEN 1 ELSE 0 END)
                         OVER (PARTITION BY player_id ORDER BY match_date) AS loss_grp

-- It behaves like
ROWS BETWEEN UNBOUNDED PRECEDING AND CURRENT ROW

-- Simple rolling window:
SELECT A,
      SUM(A) OVER (ORDER BY some_col ROWS BETWEEN 3 PRECEDING AND CURRENT ROW) AS rolling_sum_4
FROM your_table;
\end{lstlisting}

\begin{lstlisting}[language=SQL, linewidth=1.3\textwidth]
WITH match_grp AS (SELECT player_id, match_date, match_result,
                   SUM(CASE WHEN match_result = 'L' THEN 1 ELSE 0 END)
                         OVER (PARTITION BY player_id ORDER BY match_date) AS loss_grp
                  FROM players_results),
    win_streaks AS (SELECT player_id, loss_grp, COUNT(*) AS streak_len
                    FROM match_grp
                    WHERE match_result = 'W'
                    GROUP BY player_id, loss_grp)
SELECT player_id, MAX(streak_len) AS max_streak
FROM win_streaks
GROUP BY player_id
HAVING MAX(streak_len) = (SELECT MAX(streak_len) FROM win_streaks);
\end{lstlisting}
\end{exm}


\begin{exm}\label{exm:10547}
({\bf \color{red} ID 10547: \href{https://platform.stratascratch.com/coding/10547-actor-rating-difference-analysis?code_type=3}{Actor Rating Difference Analysis}})\\
You are given a dataset of actors and the films they have been involved in, 
including each film's release date and rating. 
For each actor, calculate the difference between the rating of their most recent film and their average rating across all previous films (the average rating excludes the most recent one).

Return a list of actors along with their average lifetime rating, the rating of their most recent film, and the difference between the two ratings. 
Round the difference calculation to 2 decimal places. 
If an actor has only one film, return zero for the difference and 
their only film’s rating for both the average and latest rating fields.
\begin{lstlisting}[language=SQL, linewidth=1.6\textwidth]
WITH rnk_film AS (SELECT actor_name, film_title, release_date, film_rating,
                         ROW_NUMBER() OVER (PARTITION BY actor_name ORDER BY release_date DESC) AS rn,
                         COUNT(*) OVER (PARTITION BY actor_name) AS film_count
                  FROM actor_rating_shift),
avg_previous AS (SELECT actor_name, AVG(film_rating) AS avg_prev_rating
                 FROM rnk_film
                 WHERE rn > 1
                 GROUP BY actor_name)
SELECT r.actor_name, 
    ROUND(CASE WHEN r.film_count = 1 THEN r.film_rating ELSE a.avg_prev_rating END, 2) AS avg_lifetime_rating,
       r.film_rating AS most_recent_rating,
    ROUND(CASE WHEN r.film_count = 1 THEN 0 ELSE r.film_rating - a.avg_prev_rating END, 2) AS rating_diff
FROM rnk_film r LEFT JOIN avg_previous a ON r.actor_name = a.actor_name
WHERE r.rn = 1
ORDER BY r.actor_name;
\end{lstlisting}
\end{exm}


\begin{exm}\label{exm:2007}
({\bf \color{red} ID 2007 : \href{https://platform.stratascratch.com/coding/2007-rank-variance-per-country?code_type=3}{Rank Variance Per Country}})\\
Compare the total number of comments made by users in each country during December 2019 and January 2020.

For each month, rank countries by their total number of comments in descending order. Countries with the same total should share the same rank, and the next rank should increase by one (without skipping numbers).

Return the names of the countries whose rank improved from December to January (that is, their rank number became smaller).\\

{\bf Recall:} When there is a tie in ranking you can handle it in 2 different ways:
\begin{itemize}[itemsep=-2pt]
\item \code{RANK()}: skips numbers (1, 1, 3)
\item  \code{DENSE\_RANK()}: no skipping (1, 1, 2)\marginnote{\idea}
\end{itemize}
\begin{lstlisting}[language=SQL, linewidth=1.35\textwidth]
WITH helper_t AS (SELECT SUM(cnt.number_of_comments) AS cmt_cnt, 
                         country.country, MONTH(cnt.created_at) AS cmt_month
                  FROM fb_comments_count cnt JOIN fb_active_users country USING (user_id)
                  WHERE created_at>='2019-12-01'AND created_at<'2020-02-01'
                  GROUP BY country.country, cmt_month),
helper_rank AS (SELECT *,  
                   DENSE_RANK() OVER (PARTITION BY cmt_month ORDER BY cmt_cnt DESC) AS rnk
                FROM helper_t),
pivoted AS (SELECT country, MAX(CASE WHEN cmt_month=12 THEN rnk END) AS dec_rank,
                            MAX(CASE WHEN cmt_month=1 THEN rnk END) AS jan_rank
            FROM helper_rank
            GROUP BY country)
SELECT country
FROM pivoted
WHERE jan_rank > dec_rank;
\end{lstlisting}
\code{pivoted} table is a table whose columns are rows of \code{helper\_rank} table to make
comparison possible.  It is doing pivoting in Python; converting tall table into wide table.
Another way would be self-join, but self-join becomes cumbersome if you have many columns to pivot. \code{MAX} in the code above is not doing nothing, we just need it to make
the \code{CASE ...} work; our table has only 1 value per month, so, there is no collapse of
rows here.
\end{exm}


\begin{exm} ({\bf \color{red} ID 10314: \href{https://platform.stratascratch.com/coding/10314-revenue-over-time?code_type=3}{Revenue Over Time}})\marginnote{\idea\idea\idea\idea}\\
Find the 3-month rolling average of total revenue from purchases given a table with users, their purchase amount, and date purchased. Do not include returns which are represented by negative purchase values. Output the year-month (YYYY-MM) and 3-month rolling average of revenue, sorted from earliest month to latest month.

A 3-month rolling average is defined by calculating the average total revenue from all user purchases for the current month and previous two months. The first two months will not be a true 3-month rolling average since we are not given data from last year. Assume each month has at least one purchase.\\

In the following, I first used \code{year\_month} as alias.
It seemed that might be a reserved word; it made trouble till it
was changed to \code{ym}. 
Also \code{DATE\_FORMAT(created\_at, '\%Y-\%m')} is a MySQL syntax. In SQLite it should be \code{strftime('\%Y-\%m', created\_at)}.
Since SQLite is not supported on stratascratch.com I used \code{DATE\_FORMAT}.
If you forgot these, you can use \code{LEFT(created\_at, 7)} in MySQL or 
\code{SUBSTR(created\_at, 1, 7)} in SQLite.\\

This solution has a bug (in general, but this question says assume each month has at least one purchase), if there is a missing month, we will not catch it.
We can create a table with all months first; see Remark~\ref{rem:createFullCalendar}
or Example~\ref{exm:10319MonthlyPerceDiff}.

\begin{lstlisting}[language=SQL, linewidth=1.4\textwidth]
WITH month_revenue AS (SELECT DATE_FORMAT(created_at,'%Y-%m') AS ym, SUM(purchase_amt) AS t_rev
                       FROM amazon_purchases
                       WHERE purchase_amt > 0
                       GROUP BY DATE_FORMAT(created_at, '%Y-%m')
                       ORDER BY ym)
SELECT ym, ROUND(AVG(t_rev) OVER (ORDER BY ym
                        ROWS BETWEEN 2 PRECEDING AND CURRENT ROW), 2) AS rolling_3_month_avg
FROM month_revenue
ORDER BY ym;
\end{lstlisting}
\end{exm}

\begin{exm} ({\bf \color{red}ID 10548: \href{https://platform.stratascratch.com/coding/10548-top-actor-ratings-by-genre?code_type=3}{Top Actor Ratings by Genre}})\\
Find the top actors based on their average movie rating within the genre they appear in most frequently.
\begin{itemize}
\item For each actor, determine their most frequent genre (i.e., the one they've appeared in the most).
\item If there is a tie in genre count, select the genre where the actor has the highest average rating.
\item If there is still a tie in both count and rating, include all tied genres for that actor.
\end{itemize}

Rank all resulting actor + genre pairs in descending order by their average movie rating.
\begin{itemize}
\item Return all pairs that fall within the top 3 ranks (not simply the top 3 rows), including ties.
\item  Do not skip rank numbers — for example, if two actors are tied at rank 1, the next rank is 2 (not 3).
\end{itemize}

\begin{lstlisting}[language=SQL, linewidth=1.4\textwidth]
with helper_ AS (SELECT actor_name, genre, count(*) as freq, AVG(movie_rating) as avg_rat
                 FROM top_actors_rating
                 GROUP BY actor_name, genre),
  max_freq_t AS (SELECT *, MAX(freq) OVER (partition by actor_name) as max_freq
                 FROM helper_),
 freq_subset AS (SELECT *
                 FROM max_freq_t
                 WHERE freq = max_freq),
maxFreq_avgRat AS (SELECT *, MAX(avg_rat) OVER (partition by actor_name) as max_avg_rat
                   FROM freq_subset),
   last_help AS (SELECT *
                 FROM maxFreq_avgRat
                 WHERE avg_rat = max_avg_rat)
SELECT *
FROM (SELECT *, DENSE_RANK() OVER (ORDER BY max_avg_rat DESC) AS rnk
      FROM last_help) t -- need t. o.w. it wont work. Nothing special about "t". Just a name.
WHERE rnk < 4
\end{lstlisting}
\end{exm}


\begin{exm} ({\bf \color{red}ID 10319: \href{https://platform.stratascratch.com/coding/10319-monthly-percentage-difference?code_type=3}{Monthly Percentage Difference}})
\label{exm:10319MonthlyPerceDiff}
\\
Given a table of purchases by date, calculate the month-over-month percentage change in revenue. The output should include the year-month date (YYYY-MM) and percentage change, rounded to the 2nd decimal point, and sorted from the beginning of the year to the end of the year.

The percentage change column will be populated from the 2nd month forward and can be calculated as 
\[ \frac{ R_t - R_{t-1}}{R_{t-1}}  \times100 \]
% $(\text{(this month's revenue - last month's revenue) / last month's revenue}) \times100 $\\
\begin{lstlisting}[language=SQL, linewidth=1.4\textwidth]
-- This solution of mine got accepted. But it has a flaw.
-- If Jan and March are present. and Feb is missing is the bug. 
-- First we should fill in the missing months.
with monthly_rev AS (SELECT DATE_FORMAT(created_at,'%Y-%m') AS ym, sum(value) as revenue
                     FROM sf_transactions
                     GROUP BY DATE_FORMAT(created_at,'%Y-%m')
                     ),
           help_ AS (SELECT *, LAG (revenue) OVER (ORDER BY ym ASC) as prev_month_rev
                     FROM monthly_rev)
SELECT ym, round(100. * (revenue - prev_month_rev)/ prev_month_rev, 2) as revenue_diff_pct
FROM help_
\end{lstlisting}

\noindent Correct answer:\marginnote{\idea \idea \idea Recursive and the whole thing}
\begin{lstlisting}[language=SQL, linewidth=1.4\textwidth]
WITH RECURSIVE months AS (SELECT DATE_FORMAT(MIN(created_at), '%Y-%m-01') AS month_start
                          FROM sf_transactions
                          UNION ALL
                          SELECT DATE_ADD(month_start, INTERVAL 1 MONTH)
                          FROM months
                          WHERE month_start < (SELECT DATE_FORMAT(MAX(created_at), '%Y-%m-01')
                                               FROM sf_transactions)
                          ),
monthly_rev AS (SELECT DATE_FORMAT(created_at,'%Y-%m-01') AS month_start, SUM(value) AS revenue
                FROM sf_transactions
                GROUP BY DATE_FORMAT(created_at,'%Y-%m-01')
                ),
full_data AS (SELECT m.month_start, COALESCE(r.revenue, 0) AS revenue
              FROM months m LEFT JOIN monthly_rev r USING (month_start) 
              ),
final AS (SELECT *, LAG(revenue) OVER (ORDER BY month_start) AS prev_month_rev
          FROM full_data)
SELECT DATE_FORMAT(month_start, '%Y-%m') AS ym,
       ROUND(100.0 * (revenue - prev_month_rev) / prev_month_rev, 2) AS revenue_diff_pct
FROM final
\end{lstlisting}
\end{exm}







\begin{exm} ({\bf \color{red}ID 2131: \href{https://platform.stratascratch.com/coding/2131-user-streaks?code_type=3}{User Streaks}})\\
Provided a table with user ID and the dates they visited the platform, find the top 3 users with the longest continuous streak of visiting the platform up to August 10, 2022. Output the user ID and the length of the streak.\\

\noindent In case of a tie, display all users with the top three longest streak lengths.\\

\noindent  Solution is below. I, of course, applied the trick I learned in \ref{ex:playersLongestStreak}. Solution is not checked by website.

In the solution below there is a line that acts like a running window, but, since no window
is specified it, by default, does running total:

\begin{lstlisting}[language=SQL, linewidth=1.6\textwidth]
SUM(CASE WHEN match_res = 'L' THEN 1 ELSE 0 END) OVER (PARTITION BY user_id ORDER BY date_visited)

-- It behaves like
ROWS BETWEEN UNBOUNDED PRECEDING AND CURRENT ROW

-- Simple rolling window:
SELECT A,
      SUM(A) OVER (ORDER BY some_column ROWS BETWEEN 3 PRECEDING AND CURRENT ROW) AS rolling_sum_4
FROM your_table;
\end{lstlisting}


\begin{lstlisting}[language=SQL, linewidth=1.6\textwidth]
WITH preAug2022 as (SELECT DISTINCT * -- date_visited, user_id
                    FROM user_streaks
                    WHERE date_visited <= '2022-08-10'
                    ORDER BY user_id, date_visited ASC),
preAug2022_lag AS (SELECT *,
                        LAG (date_visited) OVER (PARTITION BY user_id ORDER BY date_visited) as visit_lag
                   FROM preAug2022),
win_or_lose AS (SELECT *, 
                      CASE WHEN DATEDIFF(date_visited, visit_lag) = 1 THEN 'W' ELSE 'L' END as match_res
                FROM preAug2022_lag),
win_group_t AS (SELECT *,
       SUM(CASE WHEN match_res = 'L' THEN 1 ELSE 0 END) OVER (PARTITION BY user_id ORDER BY date_visited)
                           AS loss_group
                FROM win_or_lose),
streak_len AS (SELECT user_id, loss_group, COUNT(*) as streak_len
               FROM win_group_t
               WHERE match_res = 'W'
               GROUP BY user_id, loss_group),
users_maxStreak AS (SELECT user_id,  MAX(streak_len) max_streak
                    FROM streak_len
                    GROUP BY user_id),
ranked_t AS (SELECT *,  RANK() OVER (ORDER BY max_streak DESC) as rnk
             FROM users_maxStreak)
SELECT *
FROM ranked_t
WHERE rnk < 4
\end{lstlisting}
\end{exm}



\begin{exm} ({\bf \color{red}ID 2136: \href{https://platform.stratascratch.com/coding/2136-customer-tracking?code_type=3}{Customer Tracking}})\\
Given users' session logs, calculate how many hours each user was active in total across all recorded sessions.\\
Note: The session starts when \code{state=1} and ends when \code{state=0}.\\

My initial solution with flaw. The flaw is that if there is any
problem with logging costumers, then, my code wont work.
My code is assuming there is always a log in followed by a log out.
\begin{lstlisting}[language=SQL, linewidth=1.4\textwidth]
with lagged AS (SELECT *,
          LAG (timestamp) OVER (PARTITION BY cust_id ORDER BY cust_id, timestamp ASC) as lag_time
          FROM cust_tracking),
session_in_sec AS (SELECT cust_id, TIMESTAMPDIFF(SECOND, lag_time, timestamp) as session_len_sec
                   FROM lagged
                   WHERE state = 0)
SELECT cust_id, sum(session_len_sec)/3600.00 as cum_sess_HR
FROM session_in_sec
GROUP BY cust_id
\end{lstlisting}

Using ChatGPT here: If there are two consecutive logins with state = 1
followed by a log out: (any other problems?). This code NEEDS to be double checked
and most likely we can create a situation where it will break the code
\begin{lstlisting}[language=SQL, linewidth=1.4\textwidth]
WITH marked AS (SELECT *,
                      LAG(state) OVER (PARTITION BY cust_id ORDER BY timestamp) AS prev_state
                FROM cust_tracking),
sessionized AS (SELECT *,
                   SUM(CASE WHEN state = 1 AND (prev_state IS NULL OR prev_state = 0)
                    THEN 1 ELSE 0 END) OVER (PARTITION BY cust_id ORDER BY timestamp) AS session_id
                FROM marked),
paired AS (SELECT cust_id, session_id,
              MIN(CASE WHEN state = 1 THEN timestamp END) AS start_time,
              MAX(CASE WHEN state = 0 THEN timestamp END) AS end_time
           FROM sessionized
           GROUP BY cust_id, session_id)
SELECT cust_id,
       SUM(TIMESTAMPDIFF(SECOND, start_time, end_time)) / 3600.0 AS total_hours
FROM paired
WHERE end_time IS NOT NULL
GROUP BY cust_id;
\end{lstlisting}
\end{exm}



\begin{exm} ({\bf \color{red}ID 2103: \href{https://platform.stratascratch.com/coding/2103-reviewed-flags-of-top-videos?code_type=3}{Reviewed flags of top videos}})\\
For the video (or videos) that received the most user flags, how many of these flags were reviewed by YouTube? Output the video ID and the corresponding number of reviewed flags.  Ignore flags that do not have a corresponding \code{flag\_id}.


The following can be compacted.
I did a brute force. Is making it compact an advatnage?
\begin{lstlisting}[language=SQL, linewidth=1.4\textwidth]
with flag_counts AS (SELECT video_id, count(flag_id) as flag_count
                     FROM user_flags
                     WHERE flag_id is NOT NULL
                     GROUP BY video_id
                     ),
      needed_IDs AS (SELECT video_id
                     FROM flag_counts
                     WHERE flag_count = (SELECT MAX(flag_count)
                                         FROM flag_counts)
                     ),
    max_flag_IDs AS (SELECT flag_id, video_id
                     FROM user_flags
                     WHERE video_id IN (SELECT video_id FROM needed_IDs) AND flag_id IS NOT NULL
                     ),
clean_review AS (SELECT *
                 FROM flag_review
                 WHERE flag_id IN (SELECT flag_id FROM max_flag_IDs) 
                         AND reviewed_date is NOT NULL),
merged AS (SELECT cr.reviewed_date, mf.video_id
           FROM clean_review cr LEFT JOIN max_flag_IDs mf USING (flag_id)
           )
SELECT video_id, count(reviewed_date)
FROM merged
GROUP BY video_id
\end{lstlisting}

\noindent Compacted:

\begin{lstlisting}[language=SQL, linewidth=1.6\textwidth]
WITH flag_counts AS (SELECT video_id, COUNT(flag_id) AS flag_count
                    FROM user_flags
                    WHERE flag_id IS NOT NULL
                    GROUP BY video_id
                    ),
      max_videos AS (SELECT video_id
                     FROM flag_counts
                     WHERE flag_count = (SELECT MAX(flag_count) FROM flag_counts)
                    )
SELECT uf.video_id, COUNT(fr.reviewed_date) AS reviewed_flags
FROM user_flags uf INNER JOIN max_videos mv ON uf.video_id = mv.video_id
                   LEFT JOIN flag_review fr ON uf.flag_id = fr.flag_id AND fr.reviewed_date IS NOT NULL
WHERE uf.flag_id IS NOT NULL
GROUP BY uf.video_id;
\end{lstlisting}
\end{exm}

\begin{exm} ({\bf  \color{red} ID 514: \href{https://platform.stratascratch.com/coding/514-marketing-campaign-success-advanced?code_type=3}{Marketing Campaign Success}})\\
You have the \code{marketing\_campaign} table, which records in-app purchases by users. 
Users making their first in-app purchase enter a marketing campaign, where they see call-to-actions for more purchases.
 Find how many users made additional purchases due to the campaign's success.

The campaign starts one day after the first purchase. Users with only one or multiple purchases on the first day do not count, nor do users who later buy only the same products from their first day.

{\color{red}{\bf need to check this}}
\begin{lstlisting}[language=SQL, linewidth=1.4\textwidth]
WITH first_day AS (SELECT user_id, MIN(created_at) AS first_date
                   FROM marketing_campaign
                   GROUP BY user_id
                   ),
later_purchases AS (SELECT mc.user_id, mc.product_id
                    FROM marketing_campaign mc JOIN first_day fd ON mc.user_id = fd.user_id
                    WHERE mc.created_at > fd.first_date
                    ),
first_day_products AS (SELECT mc.user_id, mc.product_id
                       FROM marketing_campaign mc JOIN first_day fd ON mc.user_id = fd.user_id
                                                        AND mc.created_at = fd.first_date
                    ),
new_products AS (SELECT lp.user_id
                 FROM later_purchases lp LEFT JOIN first_day_products fp
                          ON lp.user_id = fp.user_id AND lp.product_id = fp.product_id
                 WHERE fp.product_id IS NULL
                  )
SELECT COUNT(DISTINCT np.user_id) AS users_with_new_purchases
FROM new_products np;

-----------
WITH first_day AS (
    SELECT mc.user_id, mc.created_at AS first_date
    FROM marketing_campaign mc
    WHERE mc.created_at = (
        SELECT MIN(mc2.created_at) 
        FROM marketing_campaign mc2
        WHERE mc2.user_id = mc.user_id
    )
),
later_purchases AS (
    SELECT mc.user_id, mc.product_id
    FROM marketing_campaign mc
    JOIN first_day fd ON mc.user_id = fd.user_id
    WHERE mc.created_at > fd.first_date
),
first_day_products AS (
    SELECT mc.user_id, mc.product_id
    FROM marketing_campaign mc
    JOIN first_day fd ON mc.user_id = fd.user_id
    WHERE mc.created_at = fd.first_date
),
new_products AS (
    SELECT DISTINCT lp.user_id
    FROM later_purchases lp
    LEFT JOIN first_day_products fp
           ON lp.user_id = fp.user_id AND lp.product_id = fp.product_id
    WHERE fp.product_id IS NULL
)
SELECT COUNT(DISTINCT np.user_id) AS users_with_new_purchases
FROM new_products np;
\end{lstlisting}
\end{exm}




\begin{exm} ({\bf \color{red}ID: \href{}{}})\\
\begin{lstlisting}[language=SQL, linewidth=1.4\textwidth]
\end{lstlisting}
\end{exm}
